% Оформляем как главу без номера, но с добавлением в оглавление
\chapter*{ВВЕДЕНИЕ}
\addcontentsline{toc}{chapter}{ВВЕДЕНИЕ}

Актуальность темы дипломного проекта обусловлена стремительным развитием технологий интернета вещей (IoT) и необходимостью автоматизации процессов контроля микроклимата в промышленных и бытовых помещениях. Своевременный мониторинг параметров окружающей среды позволяет предотвратить выход из строя дорогостоящего серверного оборудования и обеспечить оптимальные условия труда.

Целью данного дипломного проекта является разработка программно-аппаратного комплекса для дистанционного сбора, обработки и визуализации данных о температуре и влажности воздуха.

Для достижения поставленной цели в работе решаются следующие задачи:
\begin{itemize}
    \item проведение анализа существующих технических решений и протоколов передачи данных;
    \item выбор и обоснование элементной базы микроконтроллерного модуля;
    \item разработка принципиальной электрической схемы и топологии печатной платы;
    \item проектирование алгоритмов работы встроенного программного обеспечения;
    \item оценка экономической эффективности внедрения системы и разработка мероприятий по охране труда.
\end{itemize}

Объектом исследования является система беспроводной передачи телеметрической информации. Предметом исследования выступают методы оптимизации энергопотребления датчиков в сетях малого радиуса действия.

В первом разделе проводится обзор литературных источников и патентный поиск. Во втором разделе описывается выбор методологии проектирования. Третий и четвертый разделы посвящены непосредственной разработке аппаратной и программной частей системы.

\vspace{0.5cm}

% ОБЯЗАТЕЛЬНЫЙ АБЗАЦ ПО СТП БГУИР:
Работа выполнена автором самостоятельно. Текст пояснительной записки является оригинальным, доля авторского текста по результатам проверки в системе «Антиплагиат» составляет 85\%.

\newpage